% Companion long-record migration for Erdős #269 round 6.
% Destination label: long:actual-rationality-bridge; also carries long:jump-constraint-majorant source extract
% Not frozen R. Not a blind \input. Labels may collide with the short note;
% assign destination labels as specified in R6LongRecordMigrations.md.
% Extracted from Type B edits/relocation_extracts.tex.

\section{The actual shell orbit and its integral branch}\label{sec:actual-orbit}

We now pass from finite kernel structure to the literal infinite
$\{2,3,5\}$ shell sum.  Let $s_a$ be the total reciprocal running-height mass
of the smooth exponent triples in $[2^a,2^{a+1})$, put
\[
 T_a=\sum_{n\ge0}s_{a+n},\qquad
 X_a=\frac{\hgt(2^a)}2T_a,
\]
and let $b_a\in\{2,6,10,30\}$ and $m_a\in\N$ be respectively the ordered
block radix and ordered block digit.  These are the literal objects, not an
abstract bounded-radix model.

\begin{theorem}[actual shell-orbit dichotomy]\label{res:actual-orbit}
The series $\sum_{a\ge0}s_a$ is summable, and for every $a\ge0$,
\[
 X_{a+1}=b_aX_a-m_a.
\]
Moreover, either $X_a\in\Z$ for some $a$, or for every $a_0$ there is
$a\ge a_0$ such that
\[
 |X_a-z|\ge\frac1{31}\qquad\hbox{for every }z\in\Z.
\]
\end{theorem}

The recurrence follows by splitting the infinite tail after its first shell
and using the exact height and ordered-digit identities.  Summability follows
from the quadratic shell-count bound and geometric dyadic decay.  The last
assertion is the bounded-radix alternative applied with $2\le b_a\le30$.
It does not eliminate the first branch.  The same alternative applies to
every positive integer multiple $Y_a=BX_a$: the recurrence
$Y_{a+1}=b_aY_a-Bm_a$ keeps the radices in $\{2,6,10,30\}$, so either some
$BX_a$ is integral, and then every later term is, or
$\operatorname{dist}(BX_a,\mathbb Z)\ge 1/31$ cofinally.  A
distance-to-integer mandate that excludes the integral branch by hypothesis
therefore leaves the remaining throat untouched.

The jump word of $\{2,3,5\}$ cannot contain three consecutive factors of $2$:
between $2^e$ and $2^{e+2}$ the integer $v=\lfloor\log_3(2^e)\rfloor+1$
satisfies $2^e<3^v<2^{e+2}$, by the defining property of the integer
logarithm.  That unique-factorisation fact is the supplied-source three-power interval
lemma (\texttt{exists\_three\_power\_strictly\_between\_two\_powers} at
\path{ErdosProblems/Erdos269/JumpConstraintMajorant.lean}:28).  The geometric majorant $Q(n)=(n^2+8n+18)/9$ used only that
every multiplier is at least $2$.  Weighting the remaining product by the
stronger constraint that at least $\lfloor(k+1)/3\rfloor$ of the next $k$
letters are at least $3$, and taking the generating function
$A(z)=(1+z/2+z^2/6)/(1-z^3/12)$, yields the strictly smaller quadratic
\[
 \widetilde Q(n)=\frac{1210n^2+9130n+18847}{11979}.
\]
The moment identities and the comparison $Q-\widetilde Q>0$ are the
supplied-source moment form
(\texttt{carryMajorantQtilde\_eq\_moments} at
\path{ErdosProblems/Erdos269/JumpConstraintMajorant.lean}:102) and the
strict improvement (\texttt{carryMajorantQtilde\_lt\_Q} at line 119).  Assembling the same shell-multiplicity argument as
Theorem~\ref{res:shell} with these weights gives $X_a\le\widetilde Q(n_a)$,
where $n_a$ counts the positive $\{2,3,5\}$-powers not exceeding $2^a$.  An
exact truncated comparison, with remainder bounded by the already-recorded
checker width $a^2+6a+11$, is
\texttt{ErdosProblems/Erdos269/check\_jump\_constraint\_majorant.py}.  This
does not change the two directions of the window-escape criterion
separately.  Irrationality implies cofinal escape for every short bound
dominated by $c(B)(n+1)^2$, $\widetilde Q$ included.  The reverse
implication uses a \emph{valid} carry-dominating cap; the registered
Lean width $90B(a+1)^2$ is such a cap, and the actual if-and-only-if is
stated for that cap.  A generic upper-quadratic bound alone is not an
equivalence: the zero cap makes least-positive-residue escape automatic.

The dyadic clearing used below is one member of a symmetric family.

\begin{theorem}[prime-power boundary clearing]\label{res:prime-power-boundary-clearing}
Let $s\in\{2,3,5\}$, let $m\ge0$, and let $x$ be a positive
$\{2,3,5\}$-smooth integer with $x<s^m$.  Then
\[
 s\,\hgt(x)\mid\hgt(s^m).
\]
\end{theorem}

\begin{proof}
Write $x=2^i3^j5^k$.  The strict inequality forces the exponent in the
$s$-channel to be at most $m-1$; monotonicity of the other two integer
logarithms already places their exponents below those of $\hgt(s^m)$.
Divisibility is therefore coordinatewise.  The exact three-channel statement
is the supplied-source checked prime-power boundary clearing theorem
(\texttt{smoothHeight\_mul\_prime\_dvd\_boundaryHeight} at
\path{ErdosProblems/Erdos269/RationalLatticeReduction.lean}:51).
\end{proof}

The $s=2$ case supplies the normaliser for the shell orbit.  The $3$- and
$5$-power boundaries are additional exact clearing families, but this paper
does not yet couple them to their own normalised tail recurrences.  Write
$h_a=\hgt(2^a)/2$, which is an integer for $a\ge1$.

\begin{theorem}[all-scale rationality lattice]\label{res:all-scale-lattice}
Every finite shell window clears at its upper normaliser: for all
$u,v\ge0$ there is $m\in\N$ such that
\[
 h_{u+v}\sum_{i=0}^{v-1}s_{u+i}=m.
\]
If $q>0$ and $T_1=p/q$ for integers $p,q$, then
\[
 qX_a\in\Z\qquad(a\ge1).
\]
Consequently there are $i<j$ and $z\in\Z$ with
\[
 X_{1+j}-X_{1+i}=z.
\]
\end{theorem}

\begin{proof}
Every denominator in the finite window divides the height at its upper
endpoint, which proves the first assertion.  Splitting $T_1$ into that finite
window and the tail at $a$, multiplying by $q h_a$, and using the clearing
identity proves $qX_a\in\Z$.  Among $q+1$ such integers two have the same
residue modulo $q$; their corresponding states differ by an integer.
\end{proof}

Both conclusions are checked as supplied-source labels:
\texttt{qsmul\_normalizedTailState\_eq\_int\_of\_value\_eq\_rat}
(integrality of $qX_a$,
\path{ErdosProblems/Erdos269/RationalLatticeReduction.lean}:196) and
\texttt{exists\_normalizedTailState\_collision\_of\_value\_eq\_rat}
(the collision, line 296).

Thus pairwise incongruence of the actual $X_a$ modulo one would prove
irrationality.  The source does not establish it: rationality predicts a
collision rather than contradicting any known orbit theorem.

The integral branch in Theorem~\ref{res:actual-orbit} has a further exact
structure.

\begin{theorem}[pinning and rigidity of the integral branch]\label{res:pinning}
For every $a\ge0$,
\[
 X_a=\frac{m_a}{b_a}+\frac{X_{a+1}}{b_a},\qquad X_a>0.
\]
If $X_a$ is integral, then $X_n$ is integral for every $n\ge a$.

For every $m,K\ge0$, the genuine state has the exact finite-depth expansion
\[
 X_m=
 \sum_{i=0}^{K-1}
   \frac{m_{m+i}}{\prod_{j=0}^{i}b_{m+j}}
 +h_mT_{m+K}.                                      \tag{14}
\]

More generally, fix $A$, a positive width function $w$, and a real orbit
$(y_n)_{n\ge A}$ satisfying $y_{n+1}=b_ny_n-m_n$.  Suppose that for every
$n\ge A$, both $y_n$ and $X_n$ lie in
\[
 \left(\frac{m_n}{b_n},\frac{m_n}{b_n}+w(n)\right],
\]
and that for every $\varepsilon>0$ there is $k_0$ such that
$w(A+k)/2^k<\varepsilon$ for all $k\ge k_0$.  Then $y_A=X_A$.
\end{theorem}

\begin{proof}
The first identity is the recurrence solved for $X_a$; positivity follows
because every shell contains the dyadic point $2^a$.  Integer coefficients in
the forward recurrence give upward closure.  Iterating the pinning identity to
depth $K$ gives the finite mixed-radix digit sum plus
$X_{m+K}/\prod_{j<K}b_{m+j}$.  The height telescope
$\hgt(2^{m+K})=\hgt(2^m)\prod_{j<K}b_{m+j}$ and
$X_{m+K}=h_{m+K}T_{m+K}$ turn that remainder exactly into $h_mT_{m+K}$,
proving~\emph{(14)}.  This is the supplied-source checked finite-depth tail telescope
(\texttt{trueNormalizedState\_eq\_telescope} at
\path{ErdosProblems/Erdos269/IntegralBranchWidth.lean}:24); no limiting
argument is hidden.

For the final assertion, the
difference of the two orbits is multiplied by $b_n\ge2$ at each step, whereas
the common window bounds its absolute value at time $A+k$ by $w(A+k)$.
Hence
\[
 2^k|y_A-X_A|\le w(A+k).
\]
The stated decay forces $y_A=X_A$.
\end{proof}

This theorem turns indefinite survival of an integral seed inside the exact
windows into equality with the genuine infinite tail.  It does not prove that
no seed survives.  The four theorem families in this section are
proved above, with the linked Lean declarations providing exact mechanical
checks of their statements.
